% ---------------------------------------------------------------
% seccoes/anexo_a.tex — Anexo A – Perguntas e respostas ao perito
% ---------------------------------------------------------------

\anexo{A}{Perguntas e respostas ao perito}

\begin{tabela}{L{\dimexpr0.1953\linewidth-2\tabcolsep\relax}L{\dimexpr0.3337\linewidth-2\tabcolsep\relax}L{\dimexpr0.4690\linewidth-2\tabcolsep\relax}}
\caption{Perguntas e respostas ao perito (sessão S2, 2 de outubro de 2026)}\label{tab:16}\\
\linhacab \cab{\textbf{Tema}} & \cab{\textbf{Pergunta (síntese)}} & \cab{\textbf{Resposta do perito (síntese)}} \\
\endfirsthead
\linhacab \cab{\textbf{Tema}} & \cab{\textbf{Pergunta (síntese)}} & \cab{\textbf{Resposta do perito (síntese)}} \\
\endhead
Falha do veículo & Que estado deve assumir o veículo em caso de falha? & «STOP imediato, sem dúvida.» \\
\rowcolor{corFundo}Relevância dos sensores & Que informação dos sensores é mais importante (velocidade, distância ao obstáculo, TTC, semáforos, meteorologia, sinais de trânsito)? & Os dados indicados são suficientes. Todos são importantes: cada sensor é necessário na situação para a qual contribui. \\
Criticidade e TTC & Como distinguir uma situação normal de uma crítica, usando o TTC para decidir entre desviar e travar? & A criticidade deve resultar de uma equação, a desenvolver pela equipa, com a velocidade atual, a distância ao objeto e a capacidade de travagem. Existe um «ponto de não-retorno» a partir do qual a colisão é inevitável. «Nunca considerem arranque e travagem instantâneos.» \\
\rowcolor{corFundo}Número de leituras & Quantas leituras considerar antes de decidir? & Depende do resultado da equação. Numa situação crítica, mas ainda passível de manobra, decide-se à primeira leitura (por exemplo, mudança de faixa). Numa situação normal, pode fazer-se uma segunda leitura para excluir um falso positivo. \\
Prioridades e desempate & Que níveis de prioridade definir e como desempatar situações com a mesma prioridade? & Quatro níveis: baixo, médio, alto e crítico. Para desempatar, escolher aleatoriamente ou, se as opções estiverem ordenadas numa lista, escolher a primeira. \\
\rowcolor{corFundo}Memória de trabalho & Durante quanto tempo guardar factos como a presença de um peão ou de um obstáculo? & Quando o objeto sai do campo de deteção passa a ser um evento passado e pode ser esquecido. Uma alteração por sinal de trânsito é uma mudança de estado, não uma posição a memorizar. \\
Dados contraditórios & Como deve o veículo comportar-se quando os sensores se contradizem? & Sem uma terceira fonte, aplica-se o critério de desempate. Se o nível for alto ou crítico, «é melhor o STOP». \\
\rowcolor{corFundo}Cenários & Que cenários apresentar? & Os discutidos em S1, podendo acrescentar-se a passadeira com peão. Sugere permitir que o utilizador clique na estrada para fazer surgir um elemento fora de contexto (por exemplo, carga caída de outra faixa). \\
Explicações & Como apresentar as explicações Why e Why Not: pelas regras ou por um texto? & As regras são obrigatórias e a informação apresentada deve estar alinhada com elas. \\
\rowcolor{corFundo}Melhorias & Que melhorias sugere? & É difícil sugerir melhorias a um sistema que ainda não existe, mas considera corretos o levantamento de requisitos e o plano de trabalhos. A pergunta deve ser repetida numa fase mais avançada. \\
\end{tabela}
